\documentclass[12pt]{article}
\usepackage[utf8]{inputenc}
\usepackage[T1]{fontenc}
\usepackage{hyperref}
\title{Mbuun Language Resources for K-12 Curriculum Development}
\author{Compiled by Claude Sonnet 5.5 (claude-sonnet-5-5) for Andr\'as Kornai}
\date{}
\begin{document}
\maketitle

\section{Introduction}
Mbuun (ISO 639-3 zmp, Glottocode mpuo1241; Bantu B87) is spoken by about 165,000 to 170,000 people around Idiofa in Kwilu Province, Democratic Republic of the Congo. Other names are Mpuun, Mboon and Imbuun. Wikipedia warns that Mpuono has sometimes been wrongly conflated with Mbuun. Search found no standalone dictionary and no school materials.

\section{Primary Educational Resources}
None were found. The language is not known to be taught in schools.

\section{Linguistic and Reference Materials}
\begin{itemize}
\item Wikipedia, Mbuun language: \url{https://en.wikipedia.org/wiki/Mbuun_language}
\item Omniglot, Mbuun (kiMbuun): \url{https://www.omniglot.com/writing/mbuun.htm}
\item Omniglot chart (PDF): \url{https://www.omniglot.com/charts/mbuun.pdf}
\item Koni Muluwa and Bostoen, \emph{Lexique compar\'e des langues bantu du Kwilu} (2015), a French--English comparative lexicon of 21 languages, with Mbuun among them: \url{https://www.koeppe.de/titel_print_lexique-compare-des-langues-bantu-du-kwilu-republique-democratique-du-congo}
\item Bostoen and Koni Muluwa (2014), Glottolog reference entry: \url{https://glottolog.org/resource/reference/id/300750}
\item Reflex (CNRS) bibliography for Buu: \url{https://reflex.cnrs.fr/Lexiques/webball/biblio.php?LGNOM=Buu}
\end{itemize}
The search summary also mentioned L\'eon Mundeke's 2006 \emph{El\'ements de phonologie et de morphologie de la langue mbuun (B.87)}, but I did not find a page for it.

\section{Religious and Community Resources}
None were found.

\section{Digital and Audio Resources}
None were found, apart from the Omniglot materials above.

\section{Curriculum Development Context}
Mbuun has a large speaker community, but published materials are very scarce. Scholarly work by Bostoen, Koni Muluwa and Mundeke could support orthography and lexicon development. A teaching curriculum would still need to be built from scratch.

\end{document}